\section{Experimental Setup}
\label{sec:setup}

The empirical basis of this paper is the archived artifact set under
\path{reports/publication_run/}. The main run used the repository's
comparison-runner path with the \texttt{full} campaign and \texttt{strongreject}
evaluation enabled, and the summary CSV in the publication bundle records one
20-probe campaign for each executed model.

\paragraph{Registry and providers.}
At the time of writing, \path{core/model_registry.py} defines 20 aliases across
five providers and eight model families: OpenAI GPT models, Anthropic Claude
models, Google Gemini models, Mistral-hosted models, and Together-backed
open-weight aliases spanning Llama 3.1, Gemma 2, Mixtral, Mistral Nemo, and
GPT-OSS. The executed run reached only seven of these aliases: five Anthropic
models and two Mistral models. The remaining 13 aliases appear in the combined
operational status table as blocked entries.

\paragraph{Campaign construction.}
The archived run used the framework's \texttt{full} campaign configuration,
which expands the default ATT\&CK-inspired taxonomy into 20 probes organized by
tactic category. The taxonomy spans reconnaissance, resource development,
initial access, privilege escalation, defense evasion, persistence,
exfiltration, impact, and goal hijacking. Under the stable comparison-runner
path, this yields one probe per technique, for 20 attempted probes per model.

\paragraph{Trace capture and reporting.}
Each model interaction is logged to per-model JSONL files under
\path{trajectories/publication_run/}. The corresponding structured reports are
stored under \path{reports/publication_run/per_model/}. The later publication
bundle generated from these artifacts lives under
\path{reports/publication_run/latex_quality/}, and its CSV exports supply the
summary tables used in this paper.

\paragraph{Evaluator configuration.}
Although \path{analysis/evaluators/strongreject.py} defaults to an OpenAI
\texttt{gpt-4o-mini} judge, the archived run overrode that default and used the
Claude Sonnet 4.6 alias \texttt{claude-sonnet-4-6} as the \sr\ judge. We
preserve that choice in the paper because the artifact set is the source of
truth. We do not rerun the experiments under a different judge simply to match
the code default.

\paragraph{Operational status table.}
The provider-status summary used in \cref{sec:reachability} is derived from a
combination of executed comparison results and later preflight checks.
Successful rows correspond to models that were actually run in the archived
comparison, while blocked rows correspond to aliases that failed preflight due
to quota, credential, or registry issues.
